\documentclass[11pt]{article}
\usepackage[margin=2.5cm]{geometry}
\usepackage[T1]{fontenc}
\usepackage[utf8]{inputenc}
\usepackage[slovene]{babel}
\usepackage{amsmath,amssymb,amsfonts}
\usepackage{array}
\usepackage{booktabs}
\usepackage{tabularx}
\usepackage{enumitem}

\begin{document}

\setlist[itemize]{leftmargin=1.5em, itemsep=0.3em, topsep=0.35em}
\renewcommand{\arraystretch}{1.2}

\section*{Preizkus znanja: racionalne funkcije}

\textbf{Ime in priimek:} \rule{6cm}{0.4pt}  \\
\textbf{Datum:} \rule{3cm}{0.4pt}  \hfill \textbf{Čas pisanja:} 60 minut \\
\textbf{Skupaj:} 40 točk  \hfill \textbf{Pripomočki:} pisalo in ravnilo; kalkulator ni potreben

\vspace{0.5cm}

\section*{Navodila}

\begin{itemize}
  \item Reši vse naloge in pri vsaki pokaži postopek.
  \item Rezultate zapiši natančno, kjer je to mogoče; pri besedilni nalogi dodaj enoto.
\end{itemize}
\vspace{0.2cm}

\begin{tabular}{|>{\raggedright\arraybackslash}p{2.7cm}|*{6}{c|}c|}
\hline
\textbf{Naloga} & 1 & 2 & 3 & 4 & 5 & 6 & \textbf{Skupaj} \\
\hline
\textbf{Možne točke} & 6 & 5 & 6 & 8 & 7 & 8 & \textbf{40} \\
\hline
\textbf{Pridobljene točke} &  &  &  &  &  &  &  \\
\hline
\end{tabular}

\subsection*{1. naloga}

Dana je funkcija
\[
 f(x)=\frac{x^2-4}{x^2-x-2}.
\]

\noindent\textbf{a)} Razstavi števec in imenovalec na faktorje. (2 t) \\
\noindent\textbf{b)} Določi definicijsko območje funkcije. (2 t) \\
\noindent\textbf{c)} Funkcijo poenostavi in določi koordinati morebitne luknje na grafu. (2 t)

\subsection*{2. naloga}

Reši enačbo in zapiši izločeno vrednost:
\[
 \frac{2x+1}{x-1}=3.
\]

\subsection*{3. naloga}

Reši neenačbo. Zapiši kritični vrednosti, predznake po intervalih in končno rešitev v intervalskem zapisu:
\[
 \frac{x-2}{x+3}<0.
\]

\subsection*{4. naloga}

Dana je funkcija
\[
 h(x)=\frac{3x^2+x-2}{x-1}.
\]

\noindent\textbf{a)} Določi definicijsko območje. (1 t) \\
\noindent\textbf{b)} Določi ničle in presečišče grafa z osjo $y$. (3 t) \\
\noindent\textbf{c)} Določi navpično in poševno asimptoto. (3 t) \\
\noindent\textbf{d)} Nariši okvirno skico ter označi dobljene ničle, presečišče in asimptoti. (1 t)

\subsection*{5. naloga}

Reši enačbo. Zapiši pogoje in preveri dobljeni rešitvi:
\[
 \frac{1}{x-2}+\frac{1}{x+2}=1.
\]

\subsection*{6. naloga}

Stroj A sam izdela naročilo v 6 urah. Stroj A in stroj B, ki delujeta hkrati, naročilo izdelata v 4 urah. Predpostavi, da oba stroja delata z enakomerno hitrostjo.

Koliko časa bi za izdelavo istega naročila potreboval stroj B sam? Zapiši enačbo, pokaži postopek in odgovor podaj z enoto.

\end{document}
